% =====================================================================
% 12b_captures.tex — plan des captures d'écran
% =====================================================================
\chapter{Captures à produire}
\label{ann:captures}

\section{Mode d'emploi}

Les captures sont destinées au rapport PDF et à la soutenance. Elles doivent
être lisibles, ordonnées et cohérentes avec le texte. La règle principale est
chronologique : la partie C s'exécute avant les tests destructifs de la partie
D. Si toutes les captures sont prises après le rejeu complet, quatre résultats
ne sont plus visibles : CC1, CE4, l'opération de 1\,500\,000~DA et CT4.

Deux solutions sont possibles. La première consiste à capturer la partie C
juste après \code{partie\_C\allowbreak\_requetes.sql}, avant la partie D. La
seconde consiste à utiliser les sorties conservées
\code{\_travail/\allowbreak\_C.out} et
\code{\_travail/\allowbreak\_FULL.out}. Il ne faut jamais réinventer un
résultat disparu.

\begin{attentionbox}[Chronologie des captures]
Les captures 39 à 44 doivent être prises avant la capture 45
(\code{fermeCompte}). Après la clôture, \code{nbCepargne()} du client 4 vaut 0,
\code{soldeTotal()} ne contient plus CC1 ni CE4, et les soldes de CC1 ne sont
plus disponibles.
\end{attentionbox}

\section{Préparation de l'environnement}

Avant les captures, activer les paramètres SQL*Plus et formater les colonnes
larges. Une capture doit contenir la commande et le résultat, avec les accents
et les chiffres lisibles.

\begin{lstlisting}[caption={Paramètres de capture dans SQL*Plus.},label={lst:captures-param}]
SET SERVEROUTPUT ON SIZE UNLIMITED
SET LINESIZE 200
SET PAGESIZE 100
SET DEFINE OFF
SET FEEDBACK ON
SET VERIFY OFF

COLUMN nom FORMAT A15
COLUMN prenom FORMAT A15
COLUMN num_compte FORMAT A10
COLUMN statut FORMAT A10
COLUMN type_op FORMAT A10
COLUMN type_alerte FORMAT A22
COLUMN adresse FORMAT A50
\end{lstlisting}

\section{Liste ordonnée des 66 captures}

{\footnotesize
\setlength{\tabcolsep}{3pt}
\begin{longtable}{@{}C{0.7cm}L{3.2cm}L{5.4cm}L{4.3cm}C{0.9cm}@{}}
\caption{Plan des 66 captures, dans l'ordre de production.}
\label{tab:captures}\\
\toprule
\entete{N°} & \entete{Titre et section} & \entete{Commande ou source} &
\entete{Ce qu'on doit voir} & \entete{Prio.} \\
\midrule
\endfirsthead
\toprule
\entete{N°} & \entete{Titre et section} & \entete{Commande ou source} &
\entete{Ce qu'on doit voir} & \entete{Prio.} \\
\midrule
\endhead
\bottomrule
\endfoot
\multicolumn{5}{@{}l}{\textbf{Environnement et UML}}\\
\addlinespace
01 & Connexion à la base &
\code{sqlplus}
\code{MASTER1/\allowbreak system@\allowbreak localhost:1521/\allowbreak XEPDB1}
puis
\code{SELECT USER FROM dual;} &
\code{MASTER1}, version Oracle, date. & P1 \\
02 & Arborescence des livrables &
Explorateur Windows sur \code{D:\textbackslash bdda} &
Scripts A à E, script maître, démo, guide, \code{uml/}. & P1 \\
03 & Diagramme UML &
\code{uml/banque.puml} ou \code{uml/banque.mermaid} &
Classes, héritage, \code{REF}, collections. & P1 \\
04 & UML en version texte &
\code{uml/banque.txt} &
Version lisible si le rendu graphique manque. & P2 \\
\addlinespace
\multicolumn{5}{@{}l}{\textbf{Partie A — schéma et données}}\\
\addlinespace
05 & Types métier &
\code{SELECT type\_name, typecode FROM user\_types ORDER BY type\_name;} &
18 types en partie A, 19 au final avec \code{AGENCE\_VOR\_TYPE}. & P1 \\
06 & Tables &
\code{SELECT table\_name FROM user\_tables ORDER BY table\_name;} &
12 tables, dont \code{SIGNATAIRE\_NT\_ST}. & P1 \\
07 & Contraintes &
\code{SELECT constraint\_name, constraint\_type, table\_name} et
\code{FROM user\_constraints} &
PK, FK, CHECK, UNIQUE. & P1 \\
08 & Clients &
\code{SELECT c.num\_client, c.nom, c.prenom, c.adresse.ville FROM client c;} &
6 clients. & P1 \\
09 & Téléphones (\code{VARRAY}) &
\code{SELECT c.num\_client, t.COLUMN\_VALUE FROM client c, TABLE(c.telephones) t;} &
1 à 3 téléphones, certains \code{NULL}. & P1 \\
10 & Comptes courants &
\code{SELECT cc.num\_compte, cc.solde, cc.statut,}
\code{CARDINALITY(cc.signataires) FROM cpt\_courant cc;} &
7 CC au départ, dont CC5 \code{BLOQUE}. & P1 \\
11 & Comptes d'épargne &
\code{SELECT ce.num\_compte, ce.solde, ce.taux\_interet, ce.plafond\_depot}
\code{FROM cpt\_epargne ce;} &
5 CE, dont CE2 à 3,49. & P1 \\
12 & Comptes à terme &
\code{SELECT ct.num\_compte, ct.montant\_bloque, ct.duree\_mois,}
\code{ct.date\_echeance FROM cpt\_terme ct;} &
4 CT, dont CT4 à 0. & P1 \\
13 & Opérations &
\code{SELECT id\_op, date\_op, montant, type\_op FROM operation;} &
20 opérations au départ. & P1 \\
14 & Signataires &
\code{SELECT cc.num\_compte, s.ref\_client.nom, s.droit FROM cpt\_courant cc,}
\code{TABLE(cc.signataires) s;} &
Droits \code{X}, \code{R}, \code{D}. & P1 \\
15 & Agences, conseillers, affectations &
\code{SELECT a.code, a.nom, a.ref\_responsable.nom FROM agence a;} et
\code{SELECT * FROM affectation;} &
3 agences, 4 conseillers, historique. & P2 \\
16 & Séquences &
\code{SELECT sequence\_name, last\_number FROM user\_sequences;} &
\code{SEQ\_OPERATION}, \code{SEQ\_ALERTE}, \code{SEQ\_AUDIT},
\code{SEQ\_AFFECTATION}. & P2 \\
\addlinespace
\multicolumn{5}{@{}l}{\textbf{Partie B — manipulations}}\\
\addlinespace
17 & B5a — client 99 &
\code{SELECT * FROM client WHERE num\_client = 99;} &
\code{OUKACI Ahmed}, Béjaïa, un téléphone. & P1 \\
18 & B5b — CC99 sans titulaire &
\code{SELECT num\_compte, solde, statut FROM cpt\_courant}
\code{WHERE num\_compte = 'CC99';} puis \code{UPDATE} et \code{COMMIT} &
Compte existant avant l'\code{UPDATE}, titulaire après. & P1 \\
19 & B5c — date par défaut &
\code{SELECT id\_op, date\_op, montant, type\_op FROM operation} &
\code{date\_op} = date du jour sans saisie. & P1 \\
20 & B5d — signataire \code{D} &
\code{SELECT s.ref\_client.nom, s.droit FROM cpt\_courant cc,}
\code{TABLE(cc.signataires) s WHERE cc.num\_compte = 'CC99';} &
Un signataire avec droit \code{D}. & P1 \\
21 & B5e — \code{change\_Portable} &
\code{SELECT c.num\_client, t.COLUMN\_VALUE FROM client c,}
\code{TABLE(c.telephones) t WHERE c.num\_client IN (3, 99);} &
Client 99 : 1 puis 2 téléphones ; client 3 : \code{0777-11-22-33}. & P1 \\
22 & B5f — \code{drop\_client} avant/après &
\code{SELECT COUNT(*) FROM client WHERE num\_client = 99;} &
1 puis 0 ; aucun orphelin. & P1 \\
23 & B5f — contrôle des dépendances &
\code{SELECT COUNT(*) FROM operation WHERE ref\_client = ...;} et
\code{SELECT COUNT(*) FROM cpt\_courant WHERE ref\_client = ...;} &
0 après suppression. & P2 \\
\addlinespace
\multicolumn{5}{@{}l}{\textbf{Partie C — requêtes (avant la partie D)}}\\
\addlinespace
24 & 6.10 — solde négatif & Requête 6.10 &
CC2 : $-50\,000$, titulaire BENALI, conseiller HAMDI. & P1 \\
25 & 6.15 — top 3 clients & Requête 6.15 &
Client 1 : 3\,348\,500 ; client 4 : 830\,000 ; client 2 : 300\,000. & P1 \\
26 & 6.16a — \code{VARRAY} & Requête 6.16a &
Client 1 : 3 téléphones, 2 renseignés. & P1 \\
27 & 6.18a / 6.20 — \emph{nested table} & Requêtes 6.18a et 6.20 &
CC7 : 3 signataires ; adresse complète. & P1 \\
28 & 6.22 — échéances & Requête 6.22 &
CT2 : 5 jours ; CT1 : 20 jours. & P1 \\
29 & 6.24 — opération $>1$~M & Requête 6.24 &
Opération de 1\,500\,000~DA sur CC1. & P1 \\
30 & 6.25 — suppression CT4 & Requête 6.25 puis \code{SELECT} après &
CT4 supprimé ; CT1, CT2, CT3 restants. & P1 \\
31 & 6.1 / 6.2 — formats & Requêtes 6.1 et 6.2 &
Colonnes lisibles, clients et comptes. & P2 \\
32 & 6.12 — client avec le plus de CE & Requête 6.12 &
Client 1 (2 CE). & P2 \\
33 & 6.23 — comptes dormants & Requête 6.23 &
CC5 et CC6, sans opération depuis 12 mois. & P2 \\
\addlinespace
\multicolumn{5}{@{}l}{\textbf{Partie D — programmation PL/SQL}}\\
\addlinespace
34 & D7a — téléphone client 3 & Bloc D7a &
\code{0777-11-22-33}. & P1 \\
35 & D7b — tous les titulaires & Bloc D7b &
6 clients, certains « non renseigné ». & P2 \\
36 & D7c — curseur paramétré & Bloc D7c, seuil 100\,000 puis
$-1\,000\,000$ &
CC2, CC3, CC5, CC6 ; puis \code{aucun\_compte}. & P1 \\
37 & D7d - \code{EXTEND}, \code{LIMIT}, \code{COUNT}, \code{LAST} & Bloc D7d &
Client 98 : 2 puis 3 téléphones ; 4\up{e} refusé. & P1 \\
38 & D7e — total intérêts A1 & Bloc D7e &
10\,150~DA. & P1 \\
39 & D8 — \code{nbCepargne} / \code{soldeTotal} &
\code{SELECT c.num\_client, c.nom, c.nbCepargne(), c.soldeTotal()}
\code{FROM client c;} &
2/1/1/1/0/0 et les soldes totaux. & P1 \\
40 & D8 — intérêts et frais &
\code{SELECT ce.num\_compte, ce.calculeInterets(), ce.sortieAnticipee()}
\code{FROM cpt\_epargne ce;} &
CE5 : 515\,000 / 475\,000. & P1 \\
41 & D8 — \code{passeMvt} & Bloc \code{passeMvt} + \code{ROLLBACK} &
CC7 : 185\,000 puis 195\,000, puis rollback. & P1 \\
42 & D8 — \code{virement} même agence & Bloc virement CC4 vers CC7 &
324\,000 / 186\,000, commission 0. & P1 \\
43 & D8 — \code{virement} inter-agences & Bloc virement CC4 vers CC3 &
323\,900 / 39\,000, commission 100. & P1 \\
44 & D8 — \code{virement} rollback & Bloc CC3 vers CC4, 60\,000 &
\code{ORA-20004}, CC4 reste à 325\,000. & P1 \\
45 & D8 — \code{fermeCompte} &
\code{fermeCompte('CE4')} puis \code{('CC1')} &
CE4 320\,000 ; CC1 2\,498\,500 et 5 opérations. & P1 \\
46 & D9 — refus statut & Bloc opération du titulaire (client 5) sur CC5 &
\code{ORA-20001}. & P1 \\
47 & D9 — refus droits & Bloc client 6 (non signataire) sur CC1 &
\code{ORA-20002}. & P1 \\
48 & D9 — refus plafond & Bloc retrait 210\,000 par le signataire X (client 3)
sur CC4 & \code{ORA-20003}. & P1 \\
49 & D9 — refus découvert & Bloc retrait 60\,000 par le titulaire (client 3)
sur CC3 & \code{ORA-20004}. & P1 \\
50 & D9 — succès + solde & Bloc retrait 1\,000 par le signataire D (client 2)
sur CC7 & 185\,000 puis 184\,000. & P1 \\
51 & D9 — alertes finales & \code{SELECT * FROM alerte;} après
\code{scan\_echeances} &
\code{ECHEANCE\_PROCHE} CT2 et \code{ECHEANCE\_DEPASSEE} CT3. & P1 \\
52 & D9 — audit & \code{SELECT * FROM "AUDIT";} pendant les tests 13 à 15 &
Lignes d'audit ; l'audit final est vide. & P2 \\
53 & D9 — \code{scan\_echeances} & \code{EXEC scan\_echeances;} puis
\code{SELECT * FROM alerte;} &
Alertes CT2/CT3, sans doublon. & P2 \\
\addlinespace
\multicolumn{5}{@{}l}{\textbf{Partie E — objet avancé}}\\
\addlinespace
54 & E10 — parrainage & Requête à double notation &
Clients 2/3/4 vers AMRANI ; 5 vers BENALI/AMRANI ; 6 vers CHERIF/AMRANI. & P1 \\
55 & E11 — 5 mois &
\code{SELECT r.annee, COUNT(*) FROM releves r, TABLE(r.mois\_vry) m}
\code{GROUP BY r.annee;} &
2026 : 5 mois. & P1 \\
56 & E11 — totaux par type & Requête double \code{TABLE} &
DEPOT 13/880\,000 ; RETRAIT 32/620\,000 ; VIREMENT 3/130\,000 ;
CB 2/35\,000. & P1 \\
57 & E12 — vue objet & \code{SELECT code, nom, ville FROM Agence\_VOR;} &
3 agences. & P1 \\
58 & E12 — collections conseillers &
\code{SELECT v.code, k.nom FROM Agence\_VOR v, TABLE(v.conseillers) k;} &
A1 : 1 ; A2 : 2 ; A3 : 1. & P1 \\
59 & E12 — collections comptes &
\code{SELECT v.code, c.num\_compte, c.solde FROM Agence\_VOR v,}
\code{TABLE(v.comptes) c;} &
A1 : CC4/CC7 ; A2 : CC2/CC5 ; A3 : CC3/CC6. & P1 \\
60 & E12 — mutation & Bloc \code{mutation} + requête des compteurs &
A2 2 vers 1, A3 1 vers 2, puis restauration. & P1 \\
61 & E12 — erreurs attendues &
\code{A3.mutation(1, 'A2')}, conseiller 99, agence A9 &
\code{ORA-20010}, puis deux \code{ORA-20011}. & P2 \\
\addlinespace
\multicolumn{5}{@{}l}{\textbf{Validation finale et rejeu}}\\
\addlinespace
62 & Zéro objet invalide &
\code{SELECT object\_name, object\_type, status FROM user\_objects}
\code{WHERE status <> 'VALID';} &
Aucune ligne sélectionnée. & P1 \\
63 & Nombre d'objets &
\code{SELECT object\_type, COUNT(*) FROM user\_objects GROUP BY object\_type;} &
FUNCTION 3, INDEX 21, LOB 1, PROCEDURE 6, SEQUENCE 4, TABLE 12, TRIGGER 6,
TYPE 26, TYPE BODY 6, VIEW 1. & P1 \\
64 & Rejeu du script unique &
\code{@projet\_bdda\_banque.sql} puis fin de sortie &
\code{FIN DU SCRIPT UNIQUE - PARTIES A A E}, aucune erreur de compilation. & P1 \\
65 & État final & Requêtes de comptage &
6 clients, 6 CC, 4 CE, 3 CT, 15 opérations, 9 signataires, 2 alertes,
9 affectations. & P1 \\
66 & Sortie complète & \code{\_travail/\_FULL.out} &
Toutes les sections A à E. & P2 \\
\end{longtable}
}

\section{Ordre recommandé}

\begin{enumerate}
  \item 01--04 : environnement et UML.
  \item 05--16 : partie A.
  \item 17--23 : partie B.
  \item 24--33 : partie C, avant la partie D.
  \item 34--53 : partie D.
  \item 54--61 : partie E.
  \item 62--66 : validation finale.
\end{enumerate}

Si tout doit être refait, suivre cet ordre :

\begin{enumerate}
  \item exécuter \code{partie\_A\allowbreak\_schema.sql},
        \code{partie\_B\allowbreak\_manipulations.sql} et
        \code{partie\_C\allowbreak\_requetes.sql} ;
  \item prendre les captures de la partie C ;
  \item exécuter \code{partie\_D\allowbreak\_programmation.sql} et
        \code{partie\_E\allowbreak\_objet\_avance.sql} ;
  \item prendre les captures des parties D et E ;
  \item terminer par \code{projet\_bdda\_banque.sql} pour prouver le rejeu
        complet.
\end{enumerate}

\section{Checklist des captures}

\begin{itemize}
  \item Chaque capture contient la commande et le résultat.
  \item Les accents et les chiffres sont lisibles.
  \item La partie C est capturée avant les suppressions de D.
  \item Les captures 6.24 et 6.25 sont présentes, ou justifiées par
        \code{\_C.out}.
  \item La capture « zéro objet invalide » est présente.
  \item La capture du rejeu complet est présente.
  \item Les captures sont numérotées et référencées dans le texte.
  \item Le diagramme UML est présent.
  \item Les captures d'erreurs ORA sont légendées « erreur attendue ».
  \item Les captures ne montrent pas de données personnelles réelles : le jeu
        de données est fictif.
\end{itemize}
